\documentclass[12pt]{article}

\usepackage{fontspec}
\newfontfamily\handwriting{a-casual-handwritten-pen}[Path=../assets/fonts/,Extension=.ttf,Scale=1.5]
\newfontfamily\comic{Comic Neue}

\usepackage[utf8]{inputenc}
\usepackage{amsmath}
\usepackage{amssymb}
\usepackage{enumitem}
\usepackage{xcolor}
\usepackage{soul}
\usepackage{tikz}
\usepackage[most]{tcolorbox}
\usepackage{titlesec}
\usepackage{multicol}
\usepackage{environ}

\usetikzlibrary{fit, positioning, arrows.meta}

% --- Custom Colors ---
\definecolor{primaryblue}{RGB}{42, 111, 151}
\definecolor{exbg}{RGB}{255, 240, 245}
\definecolor{rulebg}{RGB}{232, 245, 233}
\definecolor{highlight}{RGB}{255, 243, 176}
\definecolor{pinkanno}{RGB}{255, 105, 180}
\definecolor{purpleanno}{RGB}{147, 112, 219}
\definecolor{solutionbg}{RGB}{245, 245, 255}

% --- Header Styling ---
\titleformat{\section}{\color{primaryblue}\normalfont\Large\bfseries}{\thesection}{1em}{}
\titleformat{\subsection}{\color{primaryblue}\normalfont\large\bfseries}{\thesubsection}{1em}{}

% --- Friendly Boxes ---
% --- Auto-numbered Example Box ---
\newcounter{examplenum}

\newtcolorbox{friendlyexbox}[1]{
    colback=exbg,
    colframe=pinkanno!50,
    arc=5pt,
    title=\textbf{#1},
    coltitle=black,
    attach title to upper,
    after title={:\quad}
}

\newenvironment{friendlyex}{%
    \refstepcounter{examplenum}%
    \begin{friendlyexbox}{Example \theexamplenum}%
}{%
    \end{friendlyexbox}%
}

\newtcolorbox{friendlyrule}[1]{
    colback=rulebg,
    colframe=primaryblue!30,
    arc=5pt,
    fonttitle=\bfseries,
    title={#1},
    coltitle=primaryblue
}

\newcommand{\funhl}[1]{\sethlcolor{highlight}\hl{#1}}

% --- Show/Hide Solutions ---
\newif\ifshowsolutions

% Uncomment the next line to show solutions.
\showsolutionstrue

\newsavebox{\solutionmeasurebox}

\NewEnviron{solutionbox}{
\ifshowsolutions
\begin{tcolorbox}[
    colback=solutionbg,
    colframe=purpleanno!45,
    arc=5pt,
    title={Solution},
    coltitle=primaryblue,
    fonttitle=\bfseries
]
\BODY
\end{tcolorbox}
\else
\sbox{\solutionmeasurebox}{\begin{minipage}{\dimexpr\linewidth-10pt\relax}\BODY\end{minipage}}
\vspace{\dimexpr\ht\solutionmeasurebox+\dp\solutionmeasurebox+3em\relax}
\fi
}

\begin{document}
\handwriting

\begin{center}
    \Large\textbf{\color{primaryblue}Module 2 Lesson 2\\Quadratic Equations} \\
    \vspace{0.2cm}
    \hrule
\end{center}

\section*{Objectives}

\begin{friendlyrule}{In this lesson, we will learn how to}
\begin{itemize}
    \item solve quadratic equations using the zero product property;
    \item solve quadratic equations using the square root property;
    \item solve quadratic equations by completing the square;
    \item solve quadratic equations using the quadratic formula;
    \item use the discriminant to determine the number and type of solutions.
\end{itemize}
\end{friendlyrule}

A \funhl{quadratic equation} in the variable $x$ is an equation of the form

\[
ax^2+bx+c=0,
\]

where $a$, $b$, and $c$ represent real numbers and $a\neq 0$.

In this lesson, we study several ways to solve quadratic equations, which is also referred to as finding the \funhl{zeros}. The solutions, or zeros, of a quadratic equation represent the points where the graph of the corresponding quadratic function crosses the $x$-axis.

\section*{1. Zero Product Property}

\begin{friendlyrule}{Zero Product Property}
If $mn=0$, then $m=0$ or $n=0$.
\end{friendlyrule}

For example,

\[
x^2-5x+6=0
\]

factors as

\[
(x-2)(x-3)=0,
\]

so $x=2$ or $x=3$.

\begin{friendlyex}{}
Use the zero product property to solve the following quadratic equations:
\begin{itemize}
    \item[(a)] $x^2-2x-15=0$
    \item[(b)] $3x^2+7x-6=0$
\end{itemize}
\end{friendlyex}

 

\begin{solutionbox}
\textbf{(a)} Factor:

\[
x^2-2x-15=(x-5)(x+3)=0.
\]

So $x-5=0$ or $x+3=0$, giving

\[
x=5\qquad\text{or}\qquad x=-3.
\]

\textbf{(b)} We need two numbers whose product is $3\cdot(-6)=-18$ and whose sum is $7$: these are $9$ and $-2$.

\[
3x^2+9x-2x-6=0.
\]

\[
3x(x+3)-2(x+3)=0.
\]

\[
(3x-2)(x+3)=0.
\]

So $3x-2=0$ or $x+3=0$, giving

\[
x=\frac23\qquad\text{or}\qquad x=-3.
\]
\end{solutionbox}

  

\section*{2. Square Root Property}

\begin{friendlyrule}{Square Root Property}
If $x^2=k$, then $x=\pm\sqrt{k}$.
\end{friendlyrule}

For example,

\[
x^2=25\ \Longrightarrow\ x=\pm 5,\qquad\qquad (x+5)^2=8\ \Longrightarrow\ x+5=\pm\sqrt8.
\]

\section*{3. Completing the Square}

A \funhl{perfect square trinomial} can be written as the square of a binomial:

\[
x^2+10x+25\ \longrightarrow\ (x+5)^2,\qquad\qquad y^2-8y+16\ \longrightarrow\ (y-4)^2.
\]

\begin{friendlyex}{}
Solve the following quadratic equations by completing the square.
\begin{itemize}
    \item[(a)] $x^2-6x+7=0$
    \item[(b)] $3x^2+7x-8=0$
\end{itemize}
\end{friendlyex}

 

\begin{solutionbox}
\textbf{(a)}

\[
x^2-6x=-7.
\]

Add $\left(\frac{-6}{2}\right)^2=9$ to both sides:

\[
x^2-6x+9=-7+9=2.
\]

\[
(x-3)^2=2.
\]

\[
x-3=\pm\sqrt2.
\]

\[
x=3\pm\sqrt2.
\]

\textbf{(b)} Divide by $3$:

\[
x^2+\frac73 x-\frac83=0.
\]

\[
x^2+\frac73 x=\frac83.
\]

Add $\left(\frac{7}{6}\right)^2=\frac{49}{36}$ to both sides:

\[
x^2+\frac73 x+\frac{49}{36}=\frac83+\frac{49}{36}=\frac{96}{36}+\frac{49}{36}=\frac{145}{36}.
\]

\[
\left(x+\frac76\right)^2=\frac{145}{36}.
\]

\[
x+\frac76=\pm\frac{\sqrt{145}}{6}.
\]

\[
x=\frac{-7\pm\sqrt{145}}{6}.
\]
\end{solutionbox}

  

\section*{4. Quadratic Formula}

\begin{friendlyrule}{Quadratic Formula}
For the quadratic equation written in standard form $ax^2+bx+c=0$, the solutions are given by

\[
x=\frac{-b\pm\sqrt{b^2-4ac}}{2a}.
\]
\end{friendlyrule}

\begin{friendlyex}{}
Solve the following quadratic equations by using the quadratic formula.
\begin{itemize}
    \item[(a)] $3x^2+7x-8=0$
    \item[(b)] $5x^2-6x=-2$
\end{itemize}
\end{friendlyex}

 

\begin{solutionbox}
\textbf{(a)} Here $a=3$, $b=7$, $c=-8$:

\[
x=\frac{-7\pm\sqrt{7^2-4(3)(-8)}}{2(3)}=\frac{-7\pm\sqrt{49+96}}{6}=\frac{-7\pm\sqrt{145}}{6}.
\]

This matches the completing-the-square result above --- a good consistency check!

\textbf{(b)} Rewrite in standard form: $5x^2-6x+2=0$, so $a=5$, $b=-6$, $c=2$:

\[
x=\frac{6\pm\sqrt{(-6)^2-4(5)(2)}}{2(5)}=\frac{6\pm\sqrt{36-40}}{10}=\frac{6\pm\sqrt{-4}}{10}=\frac{6\pm2i}{10}.
\]

\[
x=\frac35\pm\frac15 i.
\]
\end{solutionbox}

  

\section*{5. The Discriminant}

\begin{friendlyrule}{Discriminant and Properties of Solutions}
For $ax^2+bx+c=0$, where $a$, $b$, $c$ are real numbers, the \funhl{discriminant} is $b^2-4ac$.
\begin{itemize}
    \item[(i)] $b^2-4ac=0$: one real-number solution;
    \item[(ii)] $b^2-4ac>0$: two distinct real-number solutions;
    \item[(iii)] $b^2-4ac<0$: two distinct imaginary solutions, complex conjugates.
\end{itemize}
For case (ii), if $b^2-4ac$ is a perfect square, the solutions are rational; otherwise, they are irrational.
\end{friendlyrule}

\begin{friendlyex}{}
Determine the properties of the solutions by computing the discriminant of each quadratic equation.
\begin{itemize}
    \item[(a)] $4x^2-10x+25=0$
    \item[(b)] $2x^2-5x-3=0$
    \item[(c)] $3x^2-2x=-5$
\end{itemize}
\end{friendlyex}

 

\begin{solutionbox}
\textbf{(a)} $a=4$, $b=-10$, $c=25$:

\[
b^2-4ac=100-4(4)(25)=100-400=-300<0.
\]

Two distinct imaginary solutions (complex conjugates).

\textbf{(b)} $a=2$, $b=-5$, $c=-3$:

\[
b^2-4ac=25-4(2)(-3)=25+24=49>0,
\]

and $49$ is a perfect square. Two distinct \textbf{rational} real-number solutions.

\textbf{(c)} Rewrite in standard form: $3x^2-2x+5=0$, so $a=3$, $b=-2$, $c=5$:

\[
b^2-4ac=4-4(3)(5)=4-60=-56<0.
\]

Two distinct imaginary solutions (complex conjugates).
\end{solutionbox}

  

\section*{Summary}

\begin{friendlyrule}{Key Ideas}
\begin{itemize}
    \item A quadratic equation has the form $ax^2+bx+c=0$, with $a\neq 0$.
    \item Zero product property: if $mn=0$, then $m=0$ or $n=0$ --- useful when the equation factors nicely.
    \item Square root property: if $x^2=k$, then $x=\pm\sqrt{k}$.
    \item Completing the square turns $x^2+bx$ into a perfect square trinomial by adding $\left(\frac{b}{2}\right)^2$.
    \item The quadratic formula $x=\dfrac{-b\pm\sqrt{b^2-4ac}}{2a}$ solves any quadratic equation.
    \item The discriminant $b^2-4ac$ tells you the number and type of solutions before you even solve.
\end{itemize}
\end{friendlyrule}

\end{document}
